\documentclass[11pt]{article}
\usepackage[utf8]{inputenc}
\usepackage{amsmath,amssymb}
\usepackage[margin=1in]{geometry}
\usepackage{hyperref}
\emergencystretch=3em

\title{The deterministic chart posterior ratio and the corner-vanishing regime}
\author{Claude, with Daniel Murfet}
\date{2026-09-07}

\begin{document}
\maketitle
\sloppy

\begin{abstract}
Astra~\#14 rank~2. For a fixed phase $x$ and amplitudes $y_\varphi,y_1$ with $y_{1,00}>0$ and equal
starting exponents $p$: if $y_{\varphi,00}\ne0$ then $Z[\varphi](N)/Z[1](N)\to y_{\varphi,00}/y_{1,00}$
(\texttt{chart\_posterior\_tendsto\_const}); if $y_{\varphi,00}=0$ but the constant coefficient
$B^\varphi_p\ne0$ then $(Z[\varphi]/Z[1])\log N\to B^\varphi_p/A^1_p$
(\texttt{chart\_posterior\_tendsto\_log}). The log coefficient vanishes exactly when the amplitude
vanishes at the corner (\texttt{leading\_log\_coeff\_eq\_zero\_iff}). So the posterior expectation of
an observable vanishing at the corner decays like $1/\log N$, not like a power of $N$: Astra~\#14's
correction of the naive vanishing-order heuristic, now a theorem. Zero \texttt{sorry}/\texttt{axiom}.
\end{abstract}

\section{The things formalised}
{\small
\begin{verbatim}
theorem leading_log_coeff_eq_zero_iff ... :
    canonA β h₁ h₂ k₁ k₂ (fun i j s => ampCoeff β x y (i, j) s) p = 0 ↔ y (0, 0) = 0
theorem chart_posterior_tendsto_const ... (hy₁0 : 0 < y₁ (0,0)) (hyφ0 : yφ (0,0) ≠ 0) :
    Tendsto (fun N => twoDAmp … (anaAmp (ampCoeff β x yφ) b)
                       / twoDAmp … (anaAmp (ampCoeff β x y₁) b))
      atTop (𝓝 (yφ (0,0) / y₁ (0,0)))
theorem chart_posterior_tendsto_log ... (hyφ0 : yφ (0,0) = 0)
    (hB : canonB … (ampCoeff β x yφ) … p ≠ 0) :
    Tendsto (fun N => twoDAmp … yφ … / twoDAmp … y₁ … * Real.log N) atTop
      (𝓝 (canonB … yφ … p / canonA … y₁ … p))
\end{verbatim}
}

\section{Mathematics}
By unit~167, $Z[\varphi]\sim A^\varphi_pN^{-p}\log N$ if $A^\varphi_p\ne0$ and
$Z[\varphi]\sim B^\varphi_pN^{-p}$ if $A^\varphi_p=0\ne B^\varphi_p$, while
$Z[1]\sim A^1_pN^{-p}\log N$ with $A^1_p>0$. Since $A_p=\frac{y_{00}}{k_1k_2}\int_0^\infty
s^{p-1}e^{-\beta s^2+\beta sx_{00}}ds$ with a positive integral, $A_p=0$ iff $y_{00}=0$, and the ratio
$A^\varphi_p/A^1_p$ is $y_{\varphi,00}/y_{1,00}$. The quotient calculus of unit~166 then gives the two
limits.

\section{Paper-facing meaning}
cor:empirical\_expectation's leading term $A_{i_0,j_0}(\varphi)/B_{k_0,p_0}$ with different leading
indices: at chart level the first case beyond the generic one is an observable vanishing at the
corner, and the theorem shows the leading indices differ in the log power, not in the exponent
of $N$ --- the posterior expectation decays like $1/\log N$. This corrects the tempting reading
``vanishing to order $m$ at the singularity improves the power of $N$'' (Astra~\#14: the power only
improves under coordinatewise divisibility, which raises the exponents $(h_i+1)/k_i$ to
$(h_i+a_i+1)/k_i$ and needs the face terms). The stochastic version (random phase converging in
distribution) follows the same pattern and is left for a later unit; the formulas for $B^\varphi_p$
(the face/finite-part coefficients of unit~121) remain implicit here.

\section{Lean notes}
\begin{itemize}
\item Naming a local hypothesis \texttt{h₁} shadows the exponent variable \texttt{h₁ : ℕ} and breaks
later applications (``expected ℕ''); use distinct names (\texttt{hE₁}).
\item \texttt{tendsto\_quotient\_mul\_log\_of\_eq} with $r=0$ takes the denominator equivalence at
$r+1$, which is definitionally $1$.
\item The coefficient-ratio identity is \texttt{field\_simp} after \texttt{set M := logMoment …}.
\end{itemize}

\end{document}
